%%% Local Variables:
%%% mode: latex
%%% TeX-master: t
%%% End:

% ============================================================
% 第六章：高级功能
% 本章介绍代码块、算法伪代码、附录，以及常见错误解决方法。
% ============================================================

\chapter{高级功能}\label{chap:advanced}

\section{代码块}

使用 \texttt{listings} 宏包插入带语法高亮的代码，该宏包已包含在模板中。

\subsection{行内代码}

短代码片段可用 \verb|\verb| 命令：
\begin{verbatim}
使用 \verb|\cite{key}| 引用文献。
\end{verbatim}

\subsection{代码块环境}

较长的代码使用 \texttt{lstlisting} 环境：

\begin{verbatim}
\begin{lstlisting}[language=Python, caption={数据处理脚本}]
import numpy as np
import matplotlib.pyplot as plt

data = np.loadtxt('spectrum.dat')
energy = data[:, 0]
counts = data[:, 1]

plt.plot(energy, counts)
plt.xlabel('Energy (MeV)')
plt.ylabel('Counts')
plt.savefig('spectrum.pdf')
\end{lstlisting}
\end{verbatim}

效果：

\begin{lstlisting}[language=Python, caption={能谱数据处理脚本示例}]
import numpy as np
import matplotlib.pyplot as plt

data = np.loadtxt('spectrum.dat')
energy = data[:, 0]
counts = data[:, 1]

plt.plot(energy, counts)
plt.xlabel('Energy (MeV)')
plt.ylabel('Counts')
plt.savefig('spectrum.pdf')
\end{lstlisting}

支持的语言包括：\texttt{Python}、\texttt{C}、\texttt{C++}、\texttt{Fortran}、
\texttt{Matlab}、\texttt{bash} 等。

\subsection{从文件插入代码}

如果代码较长，可以直接读取外部文件：

\begin{verbatim}
\lstinputlisting[language=C++,
    caption={主程序}]{code/main.cpp}
\end{verbatim}

\section{算法伪代码}

使用 \texttt{algorithm2e} 环境排版算法伪代码，该宏包已包含在模板中。

\begin{verbatim}
\begin{algorithm}[H]
\caption{蒙特卡洛模拟算法}
\SetAlgoLined
初始化粒子位置和动量\;
\While{未达到收敛条件}{
    随机抽样粒子初始状态\;
    \For{每个粒子}{
        跟踪粒子在介质中的传输\;
        记录能量沉积\;
    }
    统计能谱分布\;
}
输出模拟结果\;
\end{algorithm}
\end{verbatim}

效果：

\normalem
\begin{algorithm}[H]
\caption{蒙特卡洛模拟算法}
\SetAlgoLined
初始化粒子位置和动量\;
\While{未达到收敛条件}{
    随机抽样粒子初始状态\;
    \For{每个粒子}{
        跟踪粒子在介质中的传输\;
        记录能量沉积\;
    }
    统计能谱分布\;
}
输出模拟结果\;
\end{algorithm}
\ULforem

常用算法命令：
\begin{compactitem}
    \item \verb|\While{条件}{内容}|：while 循环
    \item \verb|\For{条件}{内容}|：for 循环
    \item \verb|\If{条件}{内容}|：if 语句
    \item \verb|\eIf{条件}{then}{else}|：if-else 语句
    \item \verb|\KwIn{输入}|、\verb|\KwOut{输出}|：输入输出说明
    \item 行末加 \verb|\;|：表示语句结束（显示分号）
\end{compactitem}

\section{附录}

附录适合放置篇幅较大的推导过程、原始数据、程序源码等内容，
正文中不必展开但有参考价值的材料。

在 \texttt{main.tex} 中取消附录的注释：
\begin{verbatim}
\backmatter
\appendix
\chapter{附录}

附录      % 取消此行注释以启用附录
\input{chapters/acknowledgements}
\end{verbatim}

在 \texttt{chapters/appendix.tex} 中写附录内容：
\begin{verbatim}
\chapter{附录标题}\label{app:label}

附录中的图表公式另编排序号，如 A.1, A.2...
与正文分开编号，自动处理。
\end{verbatim}

\section{物理量和单位}

核物理论文中常用的特殊写法：

\begin{compactitem}
    \item 核素符号：\verb|$^{208}$Pb| → $^{208}$Pb；\verb|$^{16}$O| → $^{16}$O
    \item 能量单位：MeV、MeV/$c^2$（质量）、MeV/$c$（动量）
    \item 截面单位：\verb|mb（毫靶）| 或 \verb|fm$^2$|；$1$ b $= 10^{-24}$ cm$^2$
    \item 运动量符号：\verb|$Q$值|、\verb|$\sigma$截面|（需要在公式中定义）
    \item 数量级：\verb|$1.5 \times 10^{-3}$| → $1.5 \times 10^{-3}$（不用 E-3）
\end{compactitem}

\section{常见编译错误及解决方法}

\subsection{undefined control sequence（未定义命令）}

\textbf{原因}：使用了未加载的宏包中的命令，或命令拼写有误。

\textbf{解决}：检查命令拼写；确认相关宏包已在 \texttt{main.tex} 或 \texttt{ustcextra.sty} 中加载。

\subsection{Overfull \\hbox（行超出版心）}

\textbf{原因}：某段文字或公式太长，\LaTeX{} 无法正常断行。

\textbf{解决}：
\begin{compactitem}
    \item 在超长英文单词中手动添加断词提示：\verb|\-|，如 \verb|super\-con\-duc\-tor|
    \item 表格列宽过窄时改用 \texttt{p\{宽度\}} 列格式
    \item 长公式使用 \texttt{align} 环境换行
\end{compactitem}

\subsection{参考文献不显示 / 显示 [?]}

\textbf{原因}：未运行 BibTeX，或 key 不匹配。

\textbf{解决}：
\begin{compactenum}
    \item 确认已按顺序运行：\texttt{xelatex → bibtex → xelatex → xelatex}
    \item 检查 \texttt{bib/refs.bib} 中的 key 与 \verb|\cite{key}| 完全一致（区分大小写）
    \item 确认 \texttt{main.tex} 中 \verb|\bibliography{bib/refs}| 路径正确
\end{compactenum}

\subsection{图片未找到}

\textbf{原因}：文件路径或文件名有误。

\textbf{解决}：
\begin{compactitem}
    \item 确认图片在 \texttt{figures/} 目录下
    \item 文件名不含中文和空格
    \item 在 \verb|\includegraphics{文件名}| 中只写文件名，不写 \texttt{figures/} 前缀（已在 \texttt{main.tex} 中通过 \verb|\graphicspath| 配置）
\end{compactitem}

\subsection{字体警告 / 找不到字体}

\textbf{原因}：系统中缺少所需中文字体。

\textbf{解决}：确保 \texttt{main.tex} 文档类选项中有 \texttt{AutoFakeFont=true}（默认已添加），
该选项在缺少字体时自动合成加粗/斜体，避免报错。

\section{有用的写作技巧}

\begin{compactitem}
    \item \textbf{快速查找符号}：访问 \texttt{detexify.kirelabs.org}，手写符号即可识别对应命令。
    \item \textbf{段落过长警告}：编译时出现 \texttt{Underfull \\hbox} 通常无害，可忽略。
    \item \textbf{保持源码整洁}：每个句子单独一行，便于 Git 追踪修改和与导师协作。
    \item \textbf{备份}：推荐用 Git 管理论文版本，每次重大修改提交一次，永远不会丢失。
\end{compactitem}

\section{本章小结}

本章介绍了代码块、算法伪代码、附录的使用方法，
以及核物理论文中常用的符号写法和常见编译错误的解决方案。
至此，本使用说明涵盖了论文写作的主要排版功能。

在正式撰写论文时，请将各章节中的说明性内容替换为你自己的研究内容。
如有疑问，可参考 \texttt{README.md} 中的常见问题部分，
或查阅《\LaTeX{} 入门》（lshort-zh-cn）文档。
